\documentclass[11pt,a4paper]{article}
\usepackage[margin=23mm,headheight=15pt]{geometry}
\usepackage{fontspec}
% A wrapper may supply font files explicitly for an isolated build.
\ifdefined\BernoulliFontPath
  \setmainfont{LiberationSerif-Regular.ttf}[Path=\BernoulliFontPath,
    BoldFont=LiberationSerif-Bold.ttf,ItalicFont=LiberationSerif-Italic.ttf,
    BoldItalicFont=LiberationSerif-BoldItalic.ttf]
  \setsansfont{LiberationSans-Regular.ttf}[Path=\BernoulliFontPath,
    BoldFont=LiberationSans-Bold.ttf,ItalicFont=LiberationSans-Italic.ttf,
    BoldItalicFont=LiberationSans-BoldItalic.ttf]
  \setmonofont{DejaVuSansMono.ttf}[Path=\BernoulliFontPath,Scale=0.82,
    BoldFont=DejaVuSansMono-Bold.ttf,ItalicFont=DejaVuSansMono-Oblique.ttf,
    BoldItalicFont=DejaVuSansMono-BoldOblique.ttf]
  \newfontfamily\namefont{NotoSerifCJKsc-Regular.otf}[Path=\BernoulliFontPath]
\else
  \setmainfont{Liberation Serif}
  \setsansfont{Liberation Sans}
  \setmonofont{DejaVu Sans Mono}[Scale=0.82]
  \newfontfamily\namefont{Noto Serif CJK SC}
\fi
\usepackage{amsmath,amssymb,amsthm,booktabs,longtable,array,enumitem}
\usepackage{xcolor,hyperref,fancyhdr,microtype}
\definecolor{ink}{HTML}{29233F}
\definecolor{accent}{HTML}{655192}
\hypersetup{colorlinks=true,linkcolor=accent,citecolor=accent,urlcolor=accent,
 pdftitle={Strict entropy concavity for weighted sums of four Bernoulli variables},
 pdfauthor={Ziyao Tang (唐梓耀); Ming Ding (丁铭); ChatGPT}}
\pagestyle{fancy}
\fancyhf{}
\fancyhead[L]{\small\sffamily Four Bernoulli sums}
\fancyhead[R]{\small\sffamily Tang, Ding and ChatGPT}
\fancyfoot[C]{\small\thepage}
\setlength{\parindent}{0pt}
\setlength{\parskip}{5pt}
\setlength{\emergencystretch}{1em}
\setlist{itemsep=2pt,topsep=4pt}
\newtheorem{theorem}{Theorem}
\newtheorem{lemma}[theorem]{Lemma}
\newtheorem{proposition}[theorem]{Proposition}
\newtheorem{corollary}[theorem]{Corollary}
\newcommand{\Bern}{\operatorname{Bernoulli}}
\newcommand{\R}{\mathbb{R}}
\newcommand{\N}{\mathbb{N}}
\newcommand{\calC}{\mathcal{C}}
\makeatletter
\let\inputtable\@@input
\makeatother
\title{\color{ink}\bfseries Strict entropy concavity for weighted sums\\
of four Bernoulli variables}
\author{Ziyao Tang {\namefont(唐梓耀)}\and
Ming Ding {\namefont(丁铭)}\and ChatGPT}
\date{26 November 2026\\[2pt]{\small Public research draft v1.0.0}}

\begin{document}
\maketitle
\begin{abstract}
Let \(B_1,\ldots,B_4\) be independent Bernoulli random variables with a
common parameter \(p\), and let \(w_1,\ldots,w_4\) be arbitrary strictly
positive real weights. We give a computer-assisted finite classification,
followed by an analytic proof, that the Shannon entropy of
\(\sum_{i=1}^4 w_iB_i\) is strictly concave in \(p\).
With natural logarithms, the sharp bound is \(H_w''(p)\le-16+6\log6\)
for \(0<p<1\), with equality exactly when the weights are all equal and
\(p=1/2\).
The classification reduces all positive real weights to 23 subset-sum
collision types. Its complete finite certificate uses only rational
arithmetic. Five elementary entropy-loss functions suffice for the
analytic part. A separate rational interval certificate checks strict
curvature on all 23 types without using those analytic estimates.
\end{abstract}

\section{Introduction and main result}
For fixed \(w=(w_1,\ldots,w_4)\in(0,\infty)^4\), set
\[
 S_p=\sum_{i=1}^4 w_iB_i,\qquad
 B_1,\ldots,B_4\overset{\mathrm{iid}}{\sim}\Bern(p),\qquad
 H_w(p)=-\sum_s \Pr(S_p=s)\log\Pr(S_p=s).
\]
Here and throughout, \(\log\) denotes the natural logarithm and
\(0\log0=0\). The weights are fixed; only their common Bernoulli
parameter varies. The sum is over the fixed finite set
\(\mathcal S_w=\{\sum_{i\in A}w_i:A\subseteq[4]\}\) of distinct
subset sums. Every probability in this set is positive for \(0<p<1\),
so \(H_w\) is smooth on that open interval.

\begin{theorem}\label{thm:main}
For every \(w\in(0,\infty)^4\), the function \(H_w''\) has its unique
maximum at \(p=1/2\), and
\[
 \boxed{H_w''(p)\le -16+6\log6<0\quad(0<p<1).}
\]
Equality in the upper bound holds exactly when \(w_1=w_2=w_3=w_4\)
and \(p=1/2\). Consequently \(H_w\) is strictly concave on \([0,1]\),
satisfies \(H_w(p)=H_w(1-p)\), and has its unique maximum at \(p=1/2\).
\end{theorem}

Neunhäuserer's Conjecture 3.1 concerns entropy concavity for recursively
overlapping finite Bernoulli distributions of arbitrary length
\cite{neunhauserer}. The law for \(n=4\) is, up to a translation and the
replacement of \(p\) by \(1-p\), the weighted sum with weights
\[
 1,\quad a_2-a_1,\quad a_3-a_2,\quad a_4-a_3,\qquad a_1=2.
\]
Theorem~\ref{thm:main} therefore proves that \(n=4\) case and also
allows arbitrary positive real weights.

The weighted formulation was discussed by Amir in the author's
MathOverflow thread \cite{amir}. That answer reports 23 four-variable
collision types, numerical verification of all 23, and analytic arguments
for seven of them. We supply an exhaustive exact classification and
analytic curvature estimates covering every type. The previously
reported count is thus supported here by a complete rational certificate.
For equal weights, the unweighted entropy-concavity theorem is
already known \cite{hillion-johnson}; our proof treats that case as
part of the same finite catalogue.

\section{A complete finite classification of positive real weights}
For \(A\subseteq[4]\), write \(s_w(A)=\sum_{i\in A}w_i\).
Partition the sixteen subsets into classes \(\calC\) with equal
weighted sums. For a class \(C\), define
\[
 N_k(C)=\#\{A\in C:|A|=k\},\qquad
 Q_C(p)=\sum_{k=0}^4 N_k(C)p^k(1-p)^{4-k}.
\]
The \emph{entropy template} is the multiset of vectors
\((N_0(C),\ldots,N_4(C))\), one for each distinct sum. Multiplicity is
retained: different sums can have the same probability polynomial.
Once this multiset is known, the locations of the support points do
not enter entropy. An \emph{equality type} is the partition of labeled
subsets into equal-sum classes, considered up to a permutation of the
four indices.

\begin{proposition}\label{prop:classification}
There are exactly 23 entropy templates for \(w\in(0,\infty)^4\).
In fact there are exactly 23 subset-sum equality types, and those
types have distinct count-vector multisets.
\end{proposition}

\begin{proof}[Proof by exhaustive rational classification]
The following reduction specifies the finite computation and why
it is complete; its implementation and all witnesses accompany this note.

\textbf{The hyperplanes.}
After cancelling common terms in \(s_w(A)=s_w(B)\), a nontrivial equality
has a mixed-sign normal \(h\in\{-1,0,1\}^4\).
Normals \(h\) and \(-h\) define the same hyperplane.
Strict positivity excludes normals with only one nonzero sign.
Thus there are
\[
 \frac{3^4-2\cdot2^4+1}{2}=25
\]
relevant hyperplanes: six equalities of one weight with one weight,
twelve of one with two, four of one with three, and three of two with two.

\textbf{All possible equality spaces.}
For an actual positive \(w\), the span \(R_w\) of its vanishing
normals has rank at most three. A basis is chosen from the 25 normals.
Row-reducing every independent set of at most three normals therefore
enumerates every possible equality space. The candidate subset counts
are \(1,25,300,2300\); canonical rational row reduction leaves the
522 spaces in Table~\ref{tab:ranks}.

\textbf{Exact positive feasibility.}
Let \(A\) be a basis matrix of a candidate row space of rank \(r\),
and put \(d=4-r\). Consider the compact polytope
\[
 P_A=\{x:Ax=0,\ {\textstyle\sum_i x_i}=1,\ x_i\ge0\}.
\]
If the normalization equation is inconsistent, this polytope is empty.
Otherwise the normalization increases the equation rank to \(r+1\).
Every vertex is obtained by adding \(d-1\) active coordinate equations
\(x_j=0\), independent modulo those equations, and solving the resulting
square system.
Enumerating every \((d-1)\)-subset of coordinates is therefore complete,
including degenerate faces. All solutions are rational.
There is a strictly positive point precisely when every coordinate is
positive at some vertex: necessity follows from the convex-hull
description of a compact polytope, and sufficiency follows by averaging
the vertices.
The 306 rejected spaces each have a certified coordinate that is zero
at all their vertices, or have no vertices.

\textbf{Excluding extra equalities.}
For a feasible space, let \(v_0,\ldots,v_{m-1}\) be primitive integer
versions of its nonnegative vertex rays. They span \(\ker A\): if
\(v>0\) is in the cone and \(y\in\ker A\), then
\(v\pm\varepsilon y\) remain in the cone for sufficiently small
\(\varepsilon>0\), so its linear span contains every such \(y\).
For \(t>0\),
\[
 w(t)=\sum_{j=0}^{m-1}t^jv_j
\]
is strictly positive. If a listed hyperplane does not contain
\(\ker A\), its value on \(w(t)\) is a nonzero polynomial of degree at
most \(m-1\). Among \(M(m-1)+1\) positive integers there is therefore
one avoiding all \(M\) such hyperplanes.
Exact substitution supplies an integer witness with precisely the
equalities forced by \(A\).
This also proves that no unrealizable equality pattern has been added.

\textbf{Recovering the templates.}
Two subset incidence vectors lie in the same class exactly when their
difference is in the row space of \(A\). This criterion is checked
against direct integer subset sums at every witness.
The 216 feasible spaces give 23 different count-vector multisets.
Permuting coordinates in all 24 ways gives 23 orbits as well, with
one different template per orbit. For every template,
\[
 \sum_{C\in\calC}(N_0(C),\ldots,N_4(C))=(1,4,6,4,1).
\]
The certificate lists every feasible witness and every infeasibility
obstruction. Its verifier reconstructs the complete row-space universe
and checks that no space is omitted or duplicated.
\end{proof}

\begin{table}[ht]
\centering
\caption{Exhaustive classification counts.}\label{tab:ranks}
\begin{tabular}{rrrr}
\toprule
Rank & Rational row spaces & Positive feasible spaces & Types modulo \(S_4\)\\
\midrule
\inputtable classification_rows.tex
\bottomrule
\end{tabular}
\end{table}

This argument covers arbitrary irrational weights as well as rational
weights, repeated weights, and multiple simultaneous collisions.
A bounded integer scan is used only to choose readable representatives;
it is not used for completeness.

\section{Entropy loss from merging atoms}
Set
\[
 q=1-p,\qquad r=pq,\qquad h(p)=-p\log p-q\log q.
\]
The four labeled bits have joint entropy \(4h(p)\).
Merging atoms of probabilities \(z_1,\ldots,z_j\) reduces entropy by
\[
 \left(\sum_i z_i\right)\log\left(\sum_i z_i\right)
 -\sum_i z_i\log z_i.
\]
Complementary subsets give complementary fibers, so they can be grouped.
The empty and full subsets are unique extreme sums and cause no loss.
A fiber of \(a\) singleton atoms, paired with its complementary fiber
of \(a\) tripleton atoms, contributes
\(ar(1-2r)\log a\).
A fiber of \(b\) doubleton atoms contributes \(br^2\log b\).

The mixed singleton--doubleton fibers and their complements contribute
\begin{align}
 F_{ab}(p)={}&pq^2\bigl[(aq+bp)\log(aq+bp)
                    -aq\log q-bp\log p\bigr]\notag\\
 &+p^2q\bigl[(ap+bq)\log(ap+bq)
                    -ap\log p-bq\log q\bigr].\label{eq:F}
\end{align}
The possible pairs are \((a,b)=(1,1),(1,2),(2,1),(1,3)\).
One singleton and its complementary tripleton in the same fiber
contribute
\[
 T(p)=r\bigl[(p^2+q^2)\log(p^2+q^2)-2p^2\log p-2q^2\log q\bigr].
\]

\textbf{Exhaustion of fiber shapes.}
If a fiber contains \(a\) singleton atoms, every doubleton in it avoids
all of those singleton indices: otherwise its weight would be strictly
larger. Consequently a mixed singleton--doubleton fiber satisfies
\(1\le b\le\binom{4-a}{2}\), giving exactly the four pairs above.
The doubleton--tripleton case follows by complementation.
If a singleton equals a tripleton, strict positivity forces the latter
to be its complement. No other subset can join that fiber: any other
singleton or doubleton avoiding that index has smaller weight, whereas
a subset containing that index and another element has larger weight.
Thus the fiber consists of precisely those two complementary atoms.
This exhausts the mixed sizes among singleton, doubleton, and tripleton
subsets, and explains why the five functions suffice.

It follows from the catalogue that every template has the exact identity
\begin{equation}\label{eq:loss}
 H_w=4h-Ar(1-2r)-Br^2
       -f_{11}F_{11}-f_{12}F_{12}-f_{21}F_{21}-f_{13}F_{13}-tT,
\end{equation}
where \(A,B\ge0\) and \(f_{11},f_{12},f_{21},f_{13},t\) are
nonnegative integers. The complete values appear in
Table~\ref{tab:budgets}. The coefficients are obtained by enumerating
the sixteen subsets and grouping equal sums. The verification program
reconstructs these coefficients and compares the generated table with
the exact source included in this manuscript.

\section{Analytic curvature bounds}
\begin{lemma}\label{lem:F}
On \(0<p<1\),
\begin{align*}
 F_{11}''&\ge-1-2\log2,&
 F_{12}''&\ge-\tfrac43-4\log3,\\
 F_{21}''&\ge-\tfrac43-2\log3,&
 F_{13}''&\ge-\tfrac32-12\log2.
\end{align*}
Each bound is attained at \(p=1/2\).
\end{lemma}

\begin{proof}
Put \(u=2p-1\), \(C=(a+b)/2\), and \(\alpha=(b-a)/(a+b)\).
Here \(|\alpha|\in\{0,1/3,1/2\}\). After substitution, the logarithms
are constants or \(\log(1\pm u)\) and \(\log(1\pm\alpha u)\).
Their Taylor series, together with their differentiated series, converge
uniformly on each compact subinterval of \(|u|<1\).
Each \(F_{ab}\) is even, so write
\(F_{ab}=\sum_{n\ge0}c_nu^{2n}\) throughout that interval.
The first coefficients obtained from \eqref{eq:F} are:
\begin{center}
\small
\begin{tabular}{c c c c}
\toprule
 &\(c_0\)&\(c_1\)&\(c_2\)\\
\midrule
\(F_{11}\)&\(\frac{\log2}{4}\)&\(-\frac{\log2}{4}-\frac18\)&\(\frac5{48}\)\\[3pt]
\(F_{12}\)&\(\frac{3\log3}{8}\)&\(-\frac{\log3}{2}-\frac16\)&\(\frac{19}{162}+\frac{\log3}{8}\)\\[3pt]
\(F_{21}\)&\(\frac{3\log3}{8}\)&\(-\frac{\log3}{4}-\frac16\)&\(\frac{25}{162}-\frac{\log3}{8}\)\\[3pt]
\(F_{13}\)&\(\log2\)&\(-\frac{3\log2}{2}-\frac3{16}\)&\(\frac{15}{128}+\frac{\log2}{2}\)\\
\bottomrule
\end{tabular}
\end{center}
All four \(c_2\) are positive. For \(F_{21}\), it suffices that
\(\log3<11/10\), since \(25/162-11/80>0\).

Here is a derivation and explicit sign check for every remaining
coefficient. Set \(F_{ab}/r=\sum_{n\ge0}g_nu^{2n}\).
The first two coefficients are
\[
 g_0=C\log(a+b),\qquad
 g_1=-\frac{ab}{a+b}-C\alpha\log(a+b).
\]
Expanding the logarithms in \eqref{eq:F} gives, for \(n\ge2\),
\begin{equation}\label{eq:gn}
 g_n=\frac{C}{2n(n-1)(2n-1)}
 \left\{\alpha^{2n-1}[n+(n-1)\alpha]-[(n-1)+n\alpha]\right\}.
\end{equation}
For completeness, the expression being expanded can also be written
\begin{align*}
 \frac{F_{ab}}r={}&C(1-\alpha u^2)\log(a+b)\\
 &+\frac C2\bigl[(1-u)(1+\alpha u)\log(1+\alpha u)
               +(1+u)(1-\alpha u)\log(1-\alpha u)\bigr]\\
 &-\frac a4\bigl[(1-u)^2\log(1-u)+(1+u)^2\log(1+u)\bigr]\\
 &-\frac b4(1-u^2)\log(1-u^2).
\end{align*}
Since \(r=(1-u^2)/4\), \(c_n=(g_n-g_{n-1})/4\) for \(n\ge3\).
Let \(D_n=n(n-1)(n-2)(2n-1)(2n-3)>0\). Direct simplification yields:
\begin{align*}
 c_n(F_{11})&=\frac18\left[
 \frac1{(n-1)(2n-3)}-\frac1{n(2n-1)}\right],\\
 c_n(F_{12})&=
 \frac{9^n(8n^2-19n+9)-(96n^3-144n^2+21n+9)}
 {8\cdot9^nD_n},\\
 c_n(F_{21})&=
 \frac{9^n(4n^2-14n+9)+(48n^3-36n^2+6n-9)}
 {8\cdot9^nD_n},\\
 c_n(F_{13})&=
 \frac{3\{4^n(4n^2-9n+4)-(12n^3-14n^2-6n+4)\}}
 {8\cdot4^nD_n}.
\end{align*}
The first is positive. For the second, when \(n\ge3\),
\[
 8n^2-19n+9\ge24,\qquad
 \frac{96n^3-144n^2+21n+9}{9^n}
 \le\frac{96n^3}{9^n}\le\frac{32}{9}<24.
\]
The last estimate follows because \(n^3/9^n\) decreases from \(n=3\).
For the third, both \(4n^2-14n+9\) and
\(48n^3-36n^2+6n-9\) are positive for \(n\ge3\).
For the fourth,
\[
 4n^2-9n+4\ge13,\qquad
 \frac{12n^3-14n^2-6n+4}{4^n}
 \le\frac{12n^3}{4^n}\le\frac{81}{16}<13,
\]
using the decrease of \(n^3/4^n\) for \(n\ge3\).
Thus \(c_n>0\) for every \(n\ge2\), in every case.

Termwise differentiation is valid on compact subintervals of \(|u|<1\).
Since \(du/dp=2\),
\[
 F_{ab}''(p)
 =8c_1+4\sum_{n\ge2}(2n)(2n-1)c_nu^{2n-2}\ge8c_1.
\]
The four values of \(8c_1\) give the stated bounds. For \(F_{13}\),
retaining the term with \(n=2\) gives the additional estimate
\begin{equation}\label{eq:quartic}
 F_{13}''(p)-F_{13}''(1/2)\ge
 48c_2(F_{13})u^2=\bigl(\tfrac{45}8+24\log2\bigr)u^2.
\end{equation}
\end{proof}

\begin{lemma}\label{lem:T}
On \(0<p<1\), \(T''(p)\ge-2=T''(1/2)\).
\end{lemma}
\begin{proof}
By symmetry, use \(u=|2p-1|\in[0,1)\), and let
\[
 O(u)=(1+u)\log(1+u)-(1-u)\log(1-u).
\]
Since \(O''(u)=-2u/(1-u^2)\le0\), \(O(0)=0\), and
\(O(1)=2\log2\) by continuity, its chord bound gives
\[
 O(u)\ge2u\log2\ge u\log\bigl(2(1+u^2)\bigr).
\]
Direct differentiation of the definition of \(T\) gives
\[
 T''(p)=-2\frac{1-u^2}{1+u^2}
 +6u\left[O(u)-u\log\bigl(2(1+u^2)\bigr)\right].
\]
The bracket is nonnegative and the first term is at least \(-2\).
The identity and the conclusion at \(u=0\) follow by continuity.
\end{proof}

\section{The sharp curvature bound}
Write \(L_2=\log2\), \(L_3=\log3\), and define
\begin{align*}
 K={}&B+f_{11}(1+2L_2)+f_{12}(\tfrac43+4L_3)\\
    &+f_{21}(\tfrac43+2L_3)+f_{13}(\tfrac32+12L_2)+2t,\\
 D={}&6A-3B-f_{13}(\tfrac{45}8+24L_2).
\end{align*}
The rational comparisons below use \(2/3<L_2<7/10\) and
\(1<L_3<11/10\). The lower bounds are the first term of
\(\log x=2\sum_{j\ge0}z^{2j+1}/(2j+1)\) for \(z=(x-1)/(x+1)\)
and \(x=2,3\). The upper bounds follow from
\(\sum_{j=0}^5(7/10)^j/j!>2\), \(\sum_{j=0}^5(11/10)^j/j!>3\),
and positivity of the omitted exponential-series terms.

\begin{proof}[Proof of Theorem~\ref{thm:main}]
Let \(r=p(1-p)\). Since
\[
 h''=-1/r,\quad [r(1-2r)]''=-6(1-4r),\quad (r^2)''=2-12r,
\]
the midpoint values in Lemmas~\ref{lem:F} and \ref{lem:T} give
\(H_w''(1/2)=-16+K\) in \eqref{eq:loss}. Subtracting the midpoint
value from \(H_w''(p)\), the \(4h\) term gives \(-4(1-4r)/r\)
and the \(A,B\) terms give \((6A-3B)(1-4r)\).
The other terms have nonpositive differences by the two lemmas,
with the stronger bound \eqref{eq:quartic} for the \(F_{13}\) term.
As \(u^2=1-4r\), this proves
\begin{equation}\label{eq:midpoint}
 H_w''(p)-H_w''(1/2)\le(1-4r)(D-4/r).
\end{equation}
In Table~\ref{tab:budgets}, the equal-weight row T01 has
\(D=30L_2-18L_3<3\). For T02,
\(D=18L_3-24L_2-45/8<-73/40\).
For T04, T05, and T08, one has \(6A=3B\) and \(f_{13}=0\).
All other rows have \(A\le2L_2\) and hence \(D\le12L_2<42/5\).
In particular \(D<16\le4/r\), so \eqref{eq:midpoint} is strictly negative
for \(p\ne1/2\).

For T01, \(K=6L_2+6L_3=6\log6>10\), whereas the rational bounds
in the last column give \(K\le99/10\) in every other row.
Thus \(H_w''(1/2)\le-16+6\log6\), with equality only for T01,
the type of equal weights. Conversely, equal weights give T01 and
attain the bound at \(p=1/2\). The bound is negative since
\(6\log6<54/5<16\).

The entropy is continuous at both endpoints because \(z\log z\to0\)
as \(z\downarrow0\), so its strict concavity in the interior extends
to \([0,1]\). Complementing all four bits sends \(S_p\) to
\(\sum_iw_i-S_p\) with the law of \(S_{1-p}\), proving symmetry.
Strict concavity and symmetry give the unique maximum at \(1/2\).
\end{proof}

\begin{corollary}
Set \(\kappa=16-6\log6\). For \(p_0,p_1,\lambda\in[0,1]\),
\begin{align*}
 H_w(\lambda p_0+(1-\lambda)p_1)\ge{}&
 \lambda H_w(p_0)+(1-\lambda)H_w(p_1)\\
 &+\frac{\kappa}{2}\lambda(1-\lambda)(p_0-p_1)^2.
\end{align*}
Equality holds exactly when \(p_0=p_1\) or \(\lambda\in\{0,1\}\).
\end{corollary}
\begin{proof}
The second derivative of \(p\mapsto H_w(p)+(\kappa/2)p^2\) is
nonpositive and vanishes at most at \(1/2\). The function extends
continuously as a concave function to \([0,1]\) and has no affine
segment of positive length. Its strict concavity yields the inequality
and the stated equality cases. For equal weights, apply the inequality
with \(p_0=1/2-\varepsilon\), \(p_1=1/2+\varepsilon\), and
\(\lambda=1/2\). As \(\varepsilon\downarrow0\), the midpoint Taylor
expansion and Theorem~\ref{thm:main} show that \(\kappa/2\) cannot be increased.
\end{proof}

\clearpage
\begin{table}[ht]
\centering
\caption{Exact entropy-loss budgets for all 23 types.
Write \(L_2=\log2\), \(L_3=\log3\), and
\(f=(f_{11},f_{12},f_{21},f_{13})\).
\(U\) is the rational upper bound for \(K\) obtained by
\(L_2<7/10,\ L_3<11/10\).}\label{tab:budgets}
\small
\setlength{\tabcolsep}{5pt}
\renewcommand{\arraystretch}{1.28}
\begin{tabular}{l c c c c c c}
\toprule
Type & Representative \(w\) & \(A\) & \(B\) & \(f\) & \(t\)& \(U\)\\
\midrule
\inputtable budget_rows.tex
\bottomrule
\end{tabular}
\end{table}

The largest rational bound for rows other than T01 is \(99/10\).
The row identifiers match the exact classification certificate.
The representatives need not be the smallest possible ones; each
is checked to have the required template.

\section{A separate exact interval check}
For a separate check of strict curvature, the supplemental program
\texttt{curvature\_certificate.py} differentiates the probability
polynomials directly, without the entropy-loss lemmas.
It uses the same exhaustive classification, with a separate calculation
of curvature from each probability polynomial.
It covers \(0\le p\le1/2\) with 136 rational closed intervals in total
across all 23 types. The other half follows from symmetry.

For each class factor
\[
 Q_C(p)=p^{a_C}q^{b_C}R_C(p),\quad
 a_C=\min\{k:N_k(C)>0\},\quad
 b_C=4-\max\{k:N_k(C)>0\}.
\]
Let \(\mathcal A=\sum_Ca_CQ_C\) and
\(\mathcal B=\sum_Cb_CQ_C\). Then
\[
 H_w=-\mathcal A\log p-\mathcal B\log q
          -\sum_C Q_C\log R_C.
\]
The verified expression is \(V(p)=-pqH_w''(p)\), written without
singular polynomial quotients:
\begin{align*}
 V={}&P+q\mathcal A''(p\log p)+p\mathcal B''(q\log q)\\
 &+pq\sum_C\left[
 Q_C''\log R_C+\frac{2Q_C'R_C'}{R_C}
 +Q_C\left(\frac{R_C''}{R_C}-\frac{(R_C')^2}{R_C^2}\right)\right],\\
 P={}&2q\mathcal A'-q\frac{\mathcal A}{p}
       -2p\mathcal B'-p\frac{\mathcal B}{q}.
\end{align*}
Both quotients in \(P\) are exact polynomial quotients.
At \(p=0\), \(V\) is understood by continuous extension; \(H_w''(0)\)
is not asserted to exist. In fact \(V(0)=4\).

All interval endpoints and all arithmetic are rational. Logarithms are
enclosed by reducing \(x=2^ky\), \(1\le y\le2\), then using
\[
 \log y=2\sum_{j=0}^{N-1}\frac{z^{2j+1}}{2j+1}+E_N,\qquad
 z=\frac{y-1}{y+1},\quad
 0\le E_N\le\frac{2z^{2N+1}}{(2N+1)(1-z^2)}.
\]
The implementation uses \(N=24\) and outward rational rounding.
For the interval meeting zero, \(p\log p\) is bounded by its monotonicity
on \([0,1/1024]\). Every certified interval has a lower bound
\(V>1/50\). This is an additional check of strict concavity;
the stronger theorem above follows from the analytic argument.

\section{Reproducibility and computational provenance}
The verification package requires only Python 3.10 or later and its
standard library. Run, without Python's optimization flag,
\begin{verbatim}
python proof/verify_all.py
\end{verbatim}
This reconstructs the finite classification universe and verifies its
certificate, reconstructs the entropy-loss identities, checks the sharp
comparison and the exact midpoint derivatives, verifies every fixed
interval certificate, and checks equality
of the three template catalogues. It also compares both displayed
classification and budget tables with their exact regenerated sources.
Finite coefficient checks in the program are transcription checks only;
the all-orders positivity proof is the analytic proof above.

The additional program \texttt{independent\_classification\_check.py}
reconstructs the classification using determinant span tests and Cramer's
rule. For each excluded space it also finds a nonzero nonnegative vector
\(y\) in that row space. Such a vector independently certifies the
exclusion: \(y\cdot w>0\) for every \(w>0\), whereas \(Aw=0\)
would imply \(y\cdot w=0\).

The main artifacts are:
\begin{itemize}
\item \texttt{classification.py}, \texttt{classification.json}:
complete real-weight coverage, witnesses, exclusions, and subset partitions;
\item \texttt{analytic\_proof.py}, \texttt{analytic\_templates.json}:
exact reconstruction of \eqref{eq:loss} and Table~\ref{tab:budgets};
\item \texttt{curvature\_certificate.py},\\
\texttt{curvature\_certificate.json}: separate rational curvature check;
\item \texttt{verify\_all.py}: coverage linkage and combined verification.
\item \texttt{manuscript\_tables.py}: deterministic generation and checking
of the classification and budget tables.
\end{itemize}

ChatGPT contributed to the derivations, code, and preparation of this
manuscript. The accompanying review note and verification records document
internal AI checks of the mathematics and computations. This public research
draft has not undergone external peer review. The finite certificates use
exact Python arithmetic; the analytic arguments appear above. No
proof-assistant kernel is used.

The conclusion concerns four positive weights with a common Bernoulli
parameter. The arbitrary-length conjecture remains a separate problem.

\paragraph{Reuse.}
This manuscript is distributed under
\href{https://creativecommons.org/licenses/by/4.0/}{CC BY 4.0}.
The verification programs and machine-readable proof certificates are
distributed under the MIT license; see the accompanying license notices.

\begin{thebibliography}{9}
\bibitem{neunhauserer}
J. Neunhäuserer,
\emph{Recursive overlap Bernoulli distributions and an entropy concavity
conjecture}, arXiv:2609.05546v2, 22 September 2026.
\url{https://arxiv.org/abs/2609.05546}.

\bibitem{amir}
J. Neunhäuserer (question), Amir (answer),
\emph{A recursive overlap construction of finite Bernoulli-sum laws
and an entropy-concavity conjecture}, MathOverflow, question 514729.
Question posted 28 August 2026; answer posted 11 September and
revised 15 September 2026. Accessed 26 November 2026.
\url{https://mathoverflow.net/questions/514729}.

\bibitem{hillion-johnson}
E. Hillion and O. Johnson,
\emph{A proof of the Shepp--Olkin entropy concavity conjecture},
Bernoulli 23(4B) (2017), 3638--3649.
\url{https://doi.org/10.3150/16-BEJ860}.
\end{thebibliography}
\end{document}
